%! Unit 1: Asking Questions with Data: Study Design — Unit Packet Cover (KEY, 2pp)
%! Page 1 is unit_cover/body.tex, shared byte-for-byte with the student cover.
%! Page 2 is the exam scoring notes for BOTH Unit 1 assessments. They live here,
%! not at the foot of a *_test_key, because the practice test is bound into the
%! STUDENT packet — shared/unit.mk merges this document into the key packet only.
\documentclass[10pt]{article}
\usepackage{saar-article}
\usepackage{saar-boxes}

\begin{document}
%! Unit 1 cover — page 1. THE SHEET ITSELF LIVES HERE.
%! Both unit_cover/main.tex and unit_cover_key/main.tex \input this file, so
%! page 1 cannot drift between the student packet and the key packet. Edit the
%! cover here, never in a wrapper.

% Full-bleed cerulean banner
\begin{tikzpicture}[remember picture, overlay]
  \fill[cerulean] (current page.north west)
    rectangle ([yshift=-1.25in]current page.north east);
\end{tikzpicture}
\vspace{-0.72in}

\begin{center}
  {\color{white}\textbf{\LARGE Statistical Analysis \& Algebraic Reasoning}}\\[4pt]
  {\color{frostmid}\large\textbf{Unit 1}}\\[6pt]
  {\color{white}\textbf{\Large Asking Questions with Data: Study Design}}\\[5pt]
  {\color{frostmid}\small\textbf{Standards: PS.DC.1 \quad PS.DC.2 \quad PS.DC.3 \quad AFDA.DA.2}}
\end{center}

\vspace{0.20in}
\namedateperiod
\vspace{0.14in}

\begin{objectivebox}
\small
Every number you will ever read about people came from a study somebody
designed. This unit is about \textbf{where data comes from} --- how a question
gets asked, who gets asked, and what the answer is then allowed to claim. By
the end you can plan a study, name what is wrong with a bad one, and tell the
difference between two sentences that look identical: \emph{is associated with}
and \emph{caused}.
\end{objectivebox}

\vspace{0.12in}

\begin{tocbox}
\small
\renewcommand{\arraystretch}{1.22}
\begin{tabularx}{\linewidth}{@{} c X l @{}}
\rowcolor{frost}
\textbf{\#} & \textbf{Lesson} & \textbf{Standards} \\
\midrule
1.0 & Unit Launch: Study Design                        & PS.DC.1a--c; AFDA.DA.2a \\
1.1 & The Statistical Cycle and Types of Data          & PS.DC.1a--c; AFDA.DA.2a \\
1.2 & Populations, Samples, Parameters, and Statistics & PS.DC.1d--e \\
1.3 & Choosing a Sample --- Four Sampling Techniques   & PS.DC.2a--b; AFDA.DA.2b \\
1.4 & Bias in Samples and Surveys                      & PS.DC.2c; AFDA.DA.2e--g \\
1.5 & Observational Studies                            & PS.DC.2d \\
1.6 & Principles of Experimental Design                & PS.DC.3a--b; AFDA.DA.2c \\
1.7 & Comparing Studies and Choosing a Method          & PS.DC.3c--e \\
1.8 & Project: Design and Conduct a Survey             & AFDA.DA.2d, h--j; PS.DC.2d \\
\midrule
    & \multicolumn{1}{r}{\textbf{Practice Test}} & \textbf{all of Unit 1} \\
\end{tabularx}
\end{tocbox}

\vspace{0.12in}

\begin{spiralbox}
\small
\begin{enumerate}[leftmargin=1.6em, itemsep=3pt, topsep=2pt]
  \item \textbf{The question decides the data.} What you ask determines whether
        the answers are categories or numbers --- before a single value exists.
  \item \textbf{A sample speaks only for the group chance could have reached.}
        A statistic describes the sample; a parameter describes the population.
  \item \textbf{Bias is a direction, not bad luck.} Honest samples scatter
        around the truth; a biased method misses the same way every time --- and
        a bigger sample repeats the error more precisely.
  \item \textbf{Who decides the groups decides the verb.} If people sorted
        themselves, you may write \emph{is associated with}. Only random
        assignment earns \emph{caused}.
\end{enumerate}
\end{spiralbox}

\vspace{0.12in}

\begin{remindbox}
\small
\textbf{The practice test at the back of this packet is yours to keep.} It
matches the real Unit 1 test in length, format, and ideas --- same four parts,
same $100$ points --- with different numbers and contexts. Work it with no
notes before you check the key, and pay attention to the questions that ask for
a \emph{direction} or a \emph{justification}: on this unit's test, naming the
right term without saying which way the error runs, or why the design earns the
verb, is worth only half the points.
\end{remindbox}


\newpage

\begin{headlinebox}{cerulean}{\color{white}\bfseries Unit 1 --- Exam Scoring Notes (Teacher Copy)}\end{headlinebox}

\vspace{4pt}
\small
\noindent
Both Unit 1 assessments are parallel forms: same blueprint, same $100$ points
(Part A $14$, Part B $16$, Part C $40$, Part D $30$), different numbers and
reshuffled vocabulary letters. Part C is item-for-item parallel, so one point
split serves both forms.
\normalsize

\begin{teachernote}[Part A \& Part B --- Answer Letters]
\scriptsize
\textbf{Practice Test.} Set 1 (1--7): \textbf{C F A B G D E}. \quad
Set 2 (8--14): \textbf{E A C G B F D}. \quad
Multiple choice (15--22): \textbf{B C A B D C A B}.

\textbf{Unit Test.} Set 1 (1--7): \textbf{A F B C G D E}. \quad
Set 2 (8--14): \textbf{C E F A D G B}. \quad
Multiple choice (15--22): \textbf{C A D B C A D B}.

\vspace{2pt}
\textbf{Why each MC answer is the right one} (the two forms test the same eight
ideas in the same order):
\textbf{15} --- the numbers are labels for groups, so the variable is
categorical; a mean of them answers nothing.
\textbf{16} --- decided by \emph{who the number describes}, never by what it
looks like; a percent can be either one.
\textbf{17} --- stratified takes some from \emph{every} group, cluster takes
everyone from a \emph{few}; both use chance.
\textbf{18} --- the test question is \emph{could that person have been chosen?}
Yes $\to$ nonresponse; no $\to$ undercoverage.
\textbf{19} --- the method did not change, so the direction and size of the
error do not change. A bigger sample repeats the error more precisely.
\textbf{20} --- nobody was assigned, so the students sorted themselves; a
lurking variable fits the data as well as the claim does.
\textbf{21} --- the control group is the baseline; random assignment is what
makes the groups alike at the start.
\textbf{22} --- the test question is \emph{can you hand this out to a person?}
No $\to$ unethical, and an unavailable experiment is a reason for a
\emph{weaker} verb, never a stronger one.
\end{teachernote}

\begin{teachernote}[Part C --- Point Split (both forms) and Answers]
\scriptsize
\textbf{Point split, item by item ($5$ each).}
\textbf{23} individuals $1$; four classifications $2$; non-statistical question
explained + repaired $2$. \quad
\textbf{24} population / sample / size $1.5$; percent and scale-up with work
$2$; parameter-vs-statistic $1$ (must say \emph{who the number describes});
census reason $0.5$. \quad
\textbf{25} fraction $1$; the shown row's work $1$; four allocations + total
$3$. \quad
\textbf{26} interval $k$ $1$; four selected numbers $2$; three techniques
$2$. \quad
\textbf{27} two percents $1.5$; two scale-ups $1.5$; biased method named, bias
type named, \textbf{direction} stated $2$ --- a correct type with no direction
earns $1$. \quad
\textbf{28} two percents + gap $1.5$; explanatory / response $1$; pooled percent
+ scale-up $1.5$; lurking variable \emph{with a direction} $1$. \quad
\textbf{29} two percents + gap $1.5$; treatment / control $1.5$; both blocks
split $2$. \quad
\textbf{30} four relative frequencies $2$; four bars drawn to those heights $2$;
scale-up $1$.

\vspace{2pt}
\textbf{Practice Test answers.}
\textbf{24} $36\%$, $540$. \quad
\textbf{25} $\frac{1}{15}$; $18$ / $16$ / $14$ / $12$; total $60$. \quad
\textbf{26} $k = 10$; $4$, $14$, $24$, $34$; cluster / convenience / simple
random. \quad
\textbf{27} $80\%$ and $45\%$; $960$ and $540$; Method 1, undercoverage,
overestimates. \quad
\textbf{28} $65\%$ and $80\%$, gap $15$; $74\%$, $666$. \quad
\textbf{29} $75\%$ and $55\%$, gap $20$; $40$/$40$ and $20$/$20$. \quad
\textbf{30} $35$ / $30$ / $25$ / $10\%$; $315$.

\textbf{Unit Test answers.}
\textbf{24} $35\%$, $420$. \quad
\textbf{25} $\frac{1}{16}$; $16$ / $14$ / $11$ / $9$; total $50$. \quad
\textbf{26} $k = 20$; $6$, $26$, $46$, $66$; stratified / convenience /
cluster. \quad
\textbf{27} $80\%$ and $30\%$; $640$ and $240$; Method 1, undercoverage,
overestimates. \quad
\textbf{28} $55\%$ and $70\%$, gap $15$; $64\%$, $384$. \quad
\textbf{29} $70\%$ and $40\%$, gap $30$; $30$/$30$ and $20$/$20$. \quad
\textbf{30} $40$ / $30$ / $20$ / $10\%$; $320$.
\end{teachernote}

\begin{teachernote}[Part D --- Rubric (both forms) and What to Watch]
\scriptsize
\textbf{31 --- the two studies ($10$).} (a) $3$: correct labels $1$, the
deciding feature \emph{who decides which group each individual is in} $2$.
(b) $4$: the experiment $1$, plus the argument that random assignment makes the
groups alike at the start so the treatment is the only difference left $3$.
\emph{Naming the experiment with no reason caps at $2$.} (c) $3$: in the
observational study the group that chose the program was already behind, so the
gap it produced was never about the program.

\vspace{2pt}
\textbf{32 --- the mail survey ($10$).} (a) $3$: response rate, percent, and
scale-up, with work. (b) $3$: nonresponse bias $1$, the gap in percentage points
$1$, the gap in households $1$. (c) $2$: recompute the same percent $1$, and say
that sample size does not remove bias $1$. (d) $2$: any change that reaches the
people who did not answer --- follow-up calls or in-person contact on a random
sample --- named together with the bias it removes.

\vspace{2pt}
\textbf{33 --- design and critique ($10$).} (a) $4$: technique $1$, justified
from what the context \emph{allows} $3$ --- a full list exists and the subgroups
are expected to differ, which is the case for stratified; accept simple random
or systematic with a justification that uses the list. (b) $3$: any of the five
methods, with a reason tied to this question. (c) $3$: names it a convenience
sample $1$, says the rest of the population had no chance of being asked $1$,
and states the one claim she may make --- about her own group only $1$.

\vspace{2pt}
\textbf{Two errors to watch for.} (i) On \textbf{33(c)}, a student who writes
``the sample is too small'' has the wrong reason --- size is not what fails,
chance is --- so cap that part at $1$. (ii) Across Part D, an answer that names
the right term but never states the \textbf{direction} of the error, or never
says \emph{why} the design earns its verb, takes half credit on that part. That
is the emphasis the cover warns students about, and it should be the emphasis
here.
\end{teachernote}

\end{document}
